% Authored. Numbers come from generated/macros.tex.

\section{Empirical Framework}
\label{sec:method}

Every quantity in this section is an accounting transformation of published data.
Nothing is estimated except the piecewise means of Section~\ref{sec:robustness},
which are reported with their sensitivity.

\subsection{Level decomposition}

The identity of Equation~\eqref{eq:identity} is applied year by year. Balances are
expressed both in million euro and as a share of contemporaneous nominal GDP, the
latter for comparability across five decades of nominal growth.

\subsection{Dynamic attribution}

Differencing the identity attributes each annual change to the subsectors:
\begin{equation}
\Delta \Bgg_t = \Delta \Bc_t + \Delta \Brl_t + \Delta \Bssf_t.
\label{eq:change}
\end{equation}

This separates the \emph{level} of a subsector balance from its \emph{contribution
to a movement}, which are distinct and which the level series alone cannot
distinguish: a subsector can be in persistent surplus and still account for a
deterioration.

Ranking episodes on nominal euro would rank them by how recent they are, since
nominal GDP in \PanelEnd\ is orders of magnitude above \PanelStart. We therefore
also report each term scaled by current-year GDP:
\begin{equation}
\frac{\Delta \Bgg_t}{GDP_t}
= \frac{\Delta \Bc_t}{GDP_t}
+ \frac{\Delta \Brl_t}{GDP_t}
+ \frac{\Delta \Bssf_t}{GDP_t}.
\label{eq:scaled}
\end{equation}

Because the four terms share one denominator, Equation~\eqref{eq:scaled} decomposes
exactly, and its residual is Equation~\eqref{eq:change}'s residual rescaled and
nothing more. It is emphatically \textbf{not} the change in the balance ratio,
$\Delta(\Bgg_t/GDP_t)$, which also moves with the denominator and does not
decompose additively. We use the scaled form for ranking and comparison and the
level form for magnitudes.

\subsection{Revenue and expenditure attribution}

For each sector-year with detailed components the balance is an identity,
$B_{i,t} = R_{i,t} - E_{i,t}$, and so is its change:
\begin{equation}
\Delta B_{i,t} = \Delta R_{i,t} - \Delta E_{i,t}.
\label{eq:revexp}
\end{equation}

Two conventions keep Equation~\eqref{eq:revexp} unambiguous. Changes are computed
only between adjacent years inside a continuous source block, so no change spans the
1995-to-2000 subsector gap. And where a term's role in the balance change is what
matters, we report its \emph{contribution}, carrying the sign with which it enters:
the contribution of expenditure is $-\Delta E$, not $\Delta E$. This matters because
a rise in expenditure reduces the balance, so placing $\Delta E$ beside $\Delta B$
invites adding two quantities of opposite sign. Contribution columns sum to $\Delta B$;
the plain change columns do not.

The attribution is also hierarchical rather than parallel. Where an episode is
attributed to a subsector, the revenue and expenditure split reported beside it is
that \emph{subsector's} own, not the aggregate's, so both halves of the table describe
the same entity.

Equations~\eqref{eq:change} and~\eqref{eq:revexp} are computed on different source
families --- the canonical balance panel and the detailed account panel. They
therefore measure the same annual change slightly differently, and we carry the
residual between them as its own column rather than reconciling it away.

\subsection{Decomposing the Social Security balance change}

For the Social Security Funds the same identity is taken one level further. The
revenue side splits into social contributions and everything else across the whole
detailed panel; the expenditure side splits into social transfers and everything else
for the modern period, where component detail exists. Writing contributions to the
balance change,
\begin{equation}
\Delta \Bssf_t
= \underbrace{\Delta C_t + \Delta R^{oth}_t}_{\text{revenue}}
\; \underbrace{- \Delta T_t - \Delta E^{oth}_t}_{\text{expenditure}},
\label{eq:ssfchange}
\end{equation}
where $C$ is social contributions, $T$ social transfers paid, and the $oth$ terms the
respective remainders. Equation~\eqref{eq:ssfchange} closes exactly by construction.
It locates an annual movement in the accounts; it does not explain it.

\subsection{Offset metric}

Writing the combined non-Social-Security balance as
$\Bnonssf_t = \Bc_t + \Brl_t$, the offset ratio is defined only when the signs make
the comparison meaningful, that is when $\Bnonssf_t < 0$ and $\Bssf_t > 0$:
\begin{equation}
O_t = \frac{\Bssf_t}{\left|\Bnonssf_t\right|}.
\label{eq:offset}
\end{equation}

$O_t = 1$ means the positive Social Security balance is exactly the size of the
negative non-Social-Security balance. The ratio compares two recorded numbers. It
is left undefined, and missing, in every year where the sign condition fails, and
it is not a counterfactual: see Section~\ref{sec:framework}.

\subsection{Persistence}

Each annual balance is classified by sign. We report, per subsector, the frequency
of each sign, the mean and median balance ratio, the longest unbroken run of each
sign, and the empirical one-year sign transition frequencies.

Two conventions follow from the 1995 splice. Magnitudes are computed inside each
regime, because a mean that averages two vintages with different levels describes
neither. Runs are \emph{not} recomputed per regime: a run is a property of the
uninterrupted series, and truncating it at a window boundary would report the
length of the window rather than the length of the run. Transition frequencies are
empirical frequencies over the transitions actually observed, not a fitted Markov
model: no standard errors and no stationarity test are claimed, and a state that
never occurs has no estimated row.

\subsection{Piecewise-constant mean shifts}

Fewer than fifty annual observations split across two regimes do not support a
flexible change-point model. The specification is deliberately restrictive: a
piecewise-constant mean with at most two breaks per regime, a minimum segment
length of five years, and exact dynamic-programming minimisation of within-segment
squared error. This is the mean-shift case of the multiple-structural-change problem
treated by \citet{baiperron1998}, and selecting the number of breaks by a Schwarz
criterion follows \citet{yao1988}.

We describe the selection rule as a \emph{BIC-style penalized criterion} rather than
as ``the BIC'', because the change-point literature uses several penalties and the
choice is not neutral. Ours charges for segment means and break locations alike:
\[
\mathrm{BIC}(k) = n \log\!\left(\frac{\mathrm{RSS}(k)}{n}\right) + (2k+1)\log n,
\]
where $k$ is the number of breaks. Charging for the locations is the conservative
choice here: it penalises an additional break twice and makes the model less willing
to claim one, which is what we want from a short sample. Zero breaks is always
admissible.

Detection runs separately inside \PanelStart--1994 and 1995--\PanelEnd, so the known
splice cannot be selected as an economic break.

Because a single selected date from eighteen or thirty-one observations is not
determined by the data, we report it with a BIC margin over the next-best break
count and with a sensitivity grid over the two tuning parameters. The results are
in Section~\ref{sec:robustness} and Appendix~\ref{sec:changepoints}, and the dates
are described as candidates throughout.
